\section{Экспериментальные результаты и анализ поведения}

Для всесторонней оценки разработанных методов было проведено тестирование в средах \texttt{antmaze-medium-navigate-v0} и \texttt{antmaze-large-navigate-v0} на 20 независимых случайных сидах (10 эпизодов на задачу, 5 задач, суммарно 1000 эпизодов на каждый метод).

\begin{table*}[t]
    \centering
    \small
    \setlength{\tabcolsep}{5.5pt}
    \caption{Сводная сравнительная таблица всех исследованных методов на AntMaze Medium и Large (20 сидов, 1000 эпизодов)}
    \label{tab:main_benchmark}
    \begin{tabular}{lcccccc}
        \toprule
        \multirow{2}{*}{\textbf{Метод планирования}} & \multicolumn{3}{c}{\textbf{AntMaze Medium}} & \multicolumn{3}{c}{\textbf{AntMaze Large}} \\
        \cmidrule(lr){2-4} \cmidrule(lr){5-7}
        & \textbf{SR (\%)} & \textbf{Шаги} & \textbf{Время (мс)} & \textbf{SR (\%)} & \textbf{Task 01 (\%)} & \textbf{Task 05 (\%)} \\
        \midrule
        1. Single-Intention Baseline & $76.0 \pm 4.2$ & 235 & 0.94 & $24.0 \pm 3.5$ & 10.0\% & 20.0\% \\
        2. Recursive Bisection & $74.2 \pm 5.1$ & 260 & 1.45 & $32.0 \pm 4.8$ & 20.0\% & 30.0\% \\
        3. Dijkstra Teacher (\texttt{high\_actor}) & $78.0 \pm 3.8$ & 210 & 0.86 & $48.0 \pm 4.2$ & 50.0\% & 50.0\% \\
        4. Distilled JAX GatedAttn [$O(1)$] & $86.0 \pm 2.9$ & 184 & \textbf{0.18} & $50.0 \pm 3.9$ & 60.0\% & 50.0\% \\
        5. Dijkstra + Single WP Translator & $\mathbf{94.0 \pm 2.1}$ & \textbf{165} & 0.56 & $44.0 \pm 4.5$ & 40.0\% & 40.0\% \\
        \textbf{6. Dijkstra + Enhanced Sequence Attn} & $88.0 \pm 3.0$ & 172 & 0.82 & $\mathbf{66.4 \pm 3.2}$ & \textbf{90.0\%} & \textbf{70.0\%} \\
        \bottomrule
    \end{tabular}
\end{table*}

\subsection{Анализ результатов на Medium лабиринте}
В компактном лабиринте Medium коридоры сравнительно короткие, и плотность поворотов высока. В этих условиях наивысшую точность продемонстрировал \texttt{Single WP Translator} (94.0\%), поскольку локальный адаптивный гейт идеально отрабатывает одиночные резкие повороты. \texttt{Distilled JAX} показал лучший баланс скорости и качества (86.0\% при латентности всего 0.18 мс).

\subsection{Триумф Sequence Reasoning на Large лабиринте}
Картина кардинально меняется при переходе к лабиринту Large, где длина коридоров увеличивается втрое, а число последовательных поворотов достигает 8–12.

На лабиринте Large базовый Single-Intention контроллер полностью деградирует (успешность падает до 24\%), а Single WP опускается до 44\%. В то же время \textbf{Enhanced Sequence Attention достигает 66.4\% (на лучших сидах до 70.0\%)} благодаря трем физическим эффектам:
\begin{enumerate}[leftmargin=*,noitemsep]
    \item \textbf{Снятие «близорукости» на длинных перегонах}: Одноточечный транслятор ориентируется строго на ближайшую точку. На высокой скорости робот не успевает погасить боковую инерцию и врезается во внутренний угол стены при входе в следующий поворот. Sequence Attention видит цепочку из 4 точек и угол кривизны $\cos \theta_k$, заблаговременно формируя плавную дугу захода.
    \item \textbf{Устранение паразитного внимания через ALiBi}: Без штрафа за расстояние трансформер «видел» далекие терминальные участки за стеной и пытался спрямлять путь. Линейный штраф ALiBi заставляет сеть отдавать 80\% веса непосредственному повороту, сохраняя глобальную ориентацию в хвосте последовательности.
    \item \textbf{Устойчивость к стохастике MuJoCo}: Скользящее окно упреждения исключает микроостановки, сохраняя равномерную походку (Gait cycle) муравья без опрокидываний.
\end{enumerate}
